% Section I (W1). Every technical sentence ends with its source tag.
\section{Introduction}\label{sec:intro}

In an emergency a clinician may need records that the normal access policy does not grant, and static access control handles this poorly in health care, where break-glass is often bolted on ad hoc~\cite{brucker2009extending}.
% src: facts-F8#6
Hospital systems add a break-the-glass override and record its use for later review~\cite{ferreira2006break}; once records are encrypted, the override must also release a key~\cite{brucker2010abe,li2013phr}.
% src: facts-F8#7, facts-F8#9, facts-F8#15

Log-gated decryption is established: in accountable and transparent decryption a key, or a share of one, is used only once the request is in a log~\cite{ryan2017accountable,kroll2018accountable,cheval2023transparency}, and Kroll et al.\ name hospital emergencies as an application~\cite{kroll2018accountable}.
% src: facts-F8#1, facts-F8#2, facts-F8#3
Break-glass for health records is established too, at the policy layer, where consent gets a lifecycle but is enforced in software~\cite{muslim2026consent}, and with keys~\cite{brucker2010abe,deoliveira2020breakglass,yang2018libac,romel2025patient}, but none of those we examined is described as gating key release on a log entry.
% src: facts-F8#17, facts-F8#9, facts-F8#13, facts-F8#14, facts-F8#18
In the works we compare (Table~\ref{tab:related}), the log-gated schemes have no patient revocation lifecycle, so none joins log-gated release to consent and revocation.
% src: facts-F8#23

This paper builds that combination and tests it.
% src: facts-F8#24
Each record is encrypted under its own data-encryption key (DEK) and moves among four lifecycle states (\texttt{GRANTED}, \texttt{REVOKED}, \texttt{EMERGENCY\_OPEN}, \texttt{EXPIRED}) by seven legal transitions, each an ECDSA-signed event.
% src: facts-F3#25, facts-F1#1, facts-F1#11, facts-F1#3, facts-F1#28
Break-glass reconstructs a per-record emergency key from three of four Shamir shares, two of which belong to the audit service, so every unlock needs the audit service.
% src: facts-F2#33, facts-F2#34, facts-F2#36
The federation is simulated: two facilities, two attribute authorities, an audit service and a key custodian on one machine, with synthetic data.
% src: facts-F4#1, facts-F5#8

Our running scenario is an unconscious patient who arrives at an out-of-network clinic (Facility~B), while the patient's records sit encrypted in a shared store and the patient's home facility is Facility~A.
The clinician reads the patient's name and date of birth from an ID card, finds the patient in the clinician portal and declares an emergency with a signed Emergency Assertion.
The audit service logs the request and only then releases its two shares; with a third share, from the clinician's key agent (the key custodian) or the patient's home authority, the facility rebuilds the record's emergency key and opens the record.
The clinician reads and searches the patient's records across facilities, and when the time window ends the record is re-keyed with no administrator action.
% src: facts-F4#1, facts-F9#7, facts-F2#33, facts-F2#34, facts-F2#36, facts-F2#50, facts-F1#47, facts-F1#55, facts-F3#51
We claim the combination, implemented, and not any of its parts:
% src: facts-F8#24
\begin{itemize}
\item \emph{One signed, logged key lifecycle.}
Patient revocation, break-glass and expiry are signed transitions of one state machine over per-record DEKs; consent is a patient-signed grant checked at read time when one is presented (status, window, signature), separate from the lifecycle state, which it does not move.
% src: facts-F1#9, facts-F1#11, facts-F1#35, facts-F1#46, facts-F1#77, facts-F1#78
Revocation and expiry rotate the DEK, the policy root, the emergency key with its shares and the search index; opening break-glass rotates nothing (Section~\ref{sec:lifecycle}).
% src: facts-F1#36, facts-F1#40, facts-F1#47
\item \emph{Transparent-decryption-style log gating for federated break-glass.}
The audit service releases its two shares only after validating a signed, one-time Emergency Assertion bound to the record and committing a hash-chained log entry (Section~\ref{sec:breakglass}).
% src: facts-F2#6, facts-F2#45, facts-F2#50, facts-F2#51, facts-F8#3
\item \emph{Emergency-scoped search.}
Search tokens derive from the emergency key, so search works only during a live emergency session, and old tokens stop matching after expiry rotation.
% src: facts-F3#51, facts-F3#53, facts-F3#54
\item \emph{Evidence.}
Negative security tests back eleven claims, C1--C11; all 35 system tests tagged to them pass in a live run on the evaluated version of the code (Sections~\ref{sec:security} and~\ref{sec:eval}), as test outcomes, not proofs; the mutation evidence for them is partial (Section~\ref{sec:security}).
% src: facts-live#12, facts-F5#22
PyShark wire checks see the audit service's database commit before the release response that carries the shares (Section~\ref{sec:traffic}).
% src: facts-F6#43, facts-F6#47, facts-F6#49
\end{itemize}
We are weaker than the transparent-decryption literature in two ways: our log is a plain hash chain held by a trusted audit service, with no proofs of presence or extension, and we give no formal model (Section~\ref{sec:limits}).
% src: facts-F8#3, facts-F8#24, facts-F2#59
